\bookchapter{第十七章：手機與其他裝置的聯動}

\chapterintro{你手裡的手機，其實是一個超級遙控器兼資料中心。它不光能打電話，還能把電腦、電視、手錶、音箱、汽車串在一起，讓這些裝置不再是孤島。這章我們就來聊聊怎麼讓它們“組隊打怪”。}

\booksection{一、與電腦協同：手機和電腦不再是兩個世界}

說實話，我以前最煩的就是手機拍了個照片，想放到電腦上修一下，還得找資料線、等傳輸。後來發現，其實根本不需要資料線，好的聯動比資料線快十倍。

\begin{itemize}
  \item \textbf{蘋果生態（AirDrop + 通用剪貼簿）}：如果你用iPhone配Mac，隔空投送是你第一個該學會的功能。在手機上點分享 隔空投送 選你的Mac，照片瞬間就過去了，不用流量不用Wi-Fi，用的是藍芽和Wi-Fi直連。還有個容易被忽略的功能叫"通用剪貼簿"——手機上覆制一段文字，在Mac上直接貼上就出來了，簡直像魔法一樣。
  \item \textbf{安卓+Windows（連線至Windows）}：微軟有個叫"連線至Windows"的App，三星手機內建，其他安卓手機去應用商店下一個就行。配對之後你可以在電腦上直接看手機通知、回簡訊、翻照片，甚至接打電話。我自己用了兩年，開會的時候不用一直低頭看手機，電腦螢幕上彈個通知就夠了。
  \item \textbf{品牌多屏協同}：華為的"多屏協同"和小米的"小米妙享"更狠——直接把手機螢幕投到電腦上，你可以用滑鼠拖拽檔案從手機到電腦，也能在電腦上刷抖音、回微信。跨裝置複製貼上、檔案拖拽都行，體驗非常絲滑。
\end{itemize}

\booksection{二、投屏到電視：全家一起看}

過年回家的時候，我堂妹想在電視上放她手機裡的畢業旅行影片。折騰了半天，最後發現其實非常簡單：

\begin{itemize}
  \item \textbf{iPhone 電視}：從右上角下拉控制中心，點那個兩個方框疊在一起的圖示（螢幕映象），找到你的電視名稱，一點就連上了。前提是電視和手機連的同一個Wi-Fi，而且電視支援AirPlay（大部分智慧電視都支援）。
  \item \textbf{安卓 電視}：各家叫法不同，華為叫"無線投屏"，小米叫"投射"，OPPO/vivo叫"多屏互動"。一般在快捷開關面板裡能找到，找不到就去設定裡搜"投屏"。
  \item \textbf{影片App直接投}：如果你只是想在電視上看愛奇藝、B站、YouTube，其實更簡單——在播放介面找一個矩形帶Wi-Fi符號的小圖示，點一下，選電視，影片就過去了。這時候手機可以幹別的，影片在電視上繼續放，互不影響。
  \item \textbf{投屏不成功怎麼辦}：90\%的情況是沒連同一個Wi-Fi。還有10\%是電視的投屏功能沒開，去電視設定裡找"多屏互動"或"AirPlay"開關。
\end{itemize}

\booksection{三、與智慧手錶/手環：健康管家長在手腕上}

我給我媽買了個手環，她一開始說“我又不運動買這個幹嘛”。結果用了兩週，她發現自己每天只睡五個多小時，開始主動調整作息了。這就是手錶/手環最實用的價值——把你看不見的健康資料變成看得見的數字。

\begin{itemize}
  \item \textbf{通知過濾器}：手機在包裡或口袋裡震個不停的時候，抬手腕看一眼——是老闆的微信還是拼多多的促銷，一秒判斷，不值得掏手機的訊息就讓它去。
  \item \textbf{健康資料同步}：手錶24小時監測心率、血氧、睡眠、步數，這些資料會自動彙總到手機上的健康App裡。久而久之你能看到自己的趨勢——這個月是不是走得太少，這周睡眠質量有沒有變好。
  \item \textbf{交通卡/門禁卡同步}：在手機上綁好的公交卡和門禁卡，大部分品牌會自動同步到手錶上。坐地鐵的時候不用掏手機，手腕一靠閘機就過去了，真的會有種"未來已來"的感覺。
  \item \textbf{找手機神器}：手錶上一般都有"找手機"功能，點一下手機就會響。我自己用這個功能的頻率遠超預期——手機隨手一放就忘在哪了。
\end{itemize}

\booksection{四、與智慧家居：手機是家裡的“總開關”}

智慧家居聽起來很玄乎，但說穿了就三件事：用手機遠端控制、設定自動化規則、語音操控。

說說我自己的幾個實際用法，你就知道這玩意兒多實用了：

\begin{itemize}
  \item \textbf{空調}：夏天回家路上提前用手機把空調開啟，進門就是涼快的。花一百多買個空調伴侶插座就能實現，不用換空調。
  \item \textbf{燈}：睡前躺床上，懶得起來關燈，手機上點一下或者喊一聲"嘿Siri/小愛同學，關燈"。
  \item \textbf{掃地機器人}：出門前手機上點一下"開始清掃"，回來地就乾淨了。還能設每天定時掃。
  \item \textbf{攝像頭}：家裡有老人小孩或者寵物的，手機上隨時能看一眼家裡的情況。
\end{itemize}

選哪個生態？一句話總結：

\begin{itemize}
  \item 用iPhone 優先選支援Apple HomeKit的裝置（貴一點但隱私保護好）。
  \item 裝置多求價效比 小米米家生態最全，燈、插座、掃地機什麼都有。
  \item 用華為手機 優先鴻蒙智聯，和華為裝置體驗最好。
\end{itemize}

\booksection{五、藍芽耳機與音箱：無線的自由}

從有線耳機換成藍芽耳機，對我來說是手機體驗升級裡最明顯的一個。再也不會從包裡掏出一團亂麻，也不用擔心線掛在門把手上了。

\begin{itemize}
  \item \textbf{配對很簡單}：耳機第一次開機一般自動進入配對模式（閃燈），手機開啟藍芽就能搜到，點一下就連上了。之後每次拿出來自動連，什麼都不用做。
  \item \textbf{同品牌有加成}：AirPods配iPhone會自動彈窗顯示電量，還能在多個蘋果裝置之間無縫切換。華為FreeBuds配華為手機、小米耳機配小米手機也都有類似的彈窗和最佳化。
  \item \textbf{藍芽音箱}：家庭聚會、戶外露營的時候，一個藍芽音箱能讓氣氛翻倍。連上之後，手機就是你的DJ臺。
  \item \textbf{注意延遲}：打遊戲的時候如果用藍芽耳機，可能會有輕微延遲（槍響了半秒才聽到槍聲）。如果對延遲敏感，要麼用有線耳機，要麼買支援"低延遲模式"的藍芽耳機。
\end{itemize}

\begin{tcbinfo}{UWB超寬頻技術——比藍芽更精準的"空間感"}
除了藍芽，新款手機裡還藏著一顆UWB（Ultra-Wideband，超寬頻）晶片。它的能力是\textbf{精確測距和空間定位}——能算出另一個裝置離你多遠、在哪個方向，精度可達釐米級。藍芽只能告訴你“在附近”，UWB能告訴你“在你左前方1.2米”。常見應用：

\begin{itemize}
  \item \textbf{AirTag/查詢精確}：iPhone查詢AirTag時顯示"7米遠、向左走"的箭頭導航，靠的就是UWB。普通藍芽只能告訴"在不在範圍內"。
  \item \textbf{數字車鑰匙}：搭載UWB的iPhone和安卓旗艦，可以當車鑰匙走近車輛自動解鎖、走遠自動鎖車，還能區分駕駛員和乘客位置，防中繼破解。
  \item \textbf{室內導航}：商場、機場裡透過UWB信標實現"過道級"導航，比GPS準得多。
  \item \textbf{支援的機型}：iPhone 11及以後機型均帶UWB晶片（U1/U2）；安卓陣營小米、三星、vivo、OPPO部分旗艦也已搭載。
\end{itemize}

\end{tcbinfo}

\begin{tcbinfo}{LE Audio與Auracast——藍芽音訊的新一代標準}
用了二十年的經典藍芽音訊（A2DP）正在被\textbf{LE Audio}（低功耗音訊）取代，這是藍芽5.2以後的新一代音訊標準：

\begin{itemize}
  \item \textbf{更省電、音質更好}：採用全新的LC3編碼，比老標準的SBC音質更好、功耗更低，耳機續航明顯延長。
  \item \textbf{雙路音訊}：左右耳塞獨立傳輸，單耳模式切換不掉線。
  \item \textbf{Auracast廣播音訊}：LE Audio的招牌功能。一個手機/電視可以同時向\textbf{無數副}耳機廣播音訊——機場候機廳、博物館、健身房、教室裡都能用一個發射源，多人用自己的耳機聽同一畫面。
  \item \textbf{現狀}：2024-2026年新上市的旗艦耳機和手機普遍支援LE Audio，老裝置需硬體升級無法軟體解鎖。買新耳機時可以留意是否支援LE Audio/Auracast標識。
\end{itemize}

\end{tcbinfo}

\booksection{六、與車載系統：開車時的安全搭檔}

開車看手機是違法的，也是要命的。但有時候確實需要導航、接電話、放音樂。這時候CarPlay（蘋果）和Android Auto（安卓）就派上用場了。

\begin{itemize}
  \item \textbf{怎麼連}：如果你的車支援（近幾年大部分新車都支援），用資料線把手機插到車上的USB口，車載螢幕會自動切換成CarPlay或Android Auto介面。部分新車支援無線連線，更省事。
  \item \textbf{能幹嘛}：導航（比車載導航好用多了）、打電話（方向盤接聽）、聽歌（手機裡的音樂App直接投到車上）、語音發微信（Siri或小愛幫忙讀訊息和回覆）。
  \item \textbf{老車也能用}：如果你的車比較老不支援，花幾百塊買個外掛的CarPlay螢幕，粘在儀表臺上，體驗差不多。
\end{itemize}

\booksection{七、聯動中的安全和隱私}

裝置聯的越多，隱私洩露的風險面就越大。幾點需要注意：

\begin{itemize}
  \item 公共場合投屏之前，先看看手機通知欄有沒有私密訊息，別讓全家看到不該看的內容。
  \item 把裝置送給別人或者賣掉之前，先解綁——去手機設定裡找到那個裝置，選"移除"或"取消配對"。
  \item 智慧家居攝像頭，買有品牌的、有加密的，別買幾十塊的雜牌。聽說過太多攝像頭被駭客直播的新聞了。
\end{itemize}

\begin{tcbpractice}{來試試：把手機螢幕投到電視上，並設定一個智慧場景}
\begin{enumerate}
  \item 確保手機和電視連同一個Wi-Fi。
  \item 在電視上找到"螢幕映象"或"AirPlay"功能並開啟。
  \item 在手機上：iPhone從控制中心點「螢幕映象」；安卓下拉找「無線投屏」。
  \item 選一張你最喜歡的照片投到電視上，讓家人看看。
  \item 如果有智慧家居裝置，試著在手機App裡建立一個"離家模式"：一鍵關燈、關空調、啟動掃地機器人。
\end{enumerate}

\end{tcbpractice}

\begin{tcbinfo}{案例 1：電池續航最佳化實戰}
\begin{itemize}
  \item \textbf{場景}：使用者小米14 續航差，半天就沒電。
  \item \textbf{步驟}：(1) 設定 電池 檢視耗電排行 (2) 識別異常後臺 App (3) 關閉 5G 切 4G (4) 限制後臺 (5) 開啟省電模式 (6) 充電時摘殼散熱。
  \item \textbf{預期結果}：續航延長 30-50\%，電池健康度資料可讀。
\end{itemize}

\end{tcbinfo}
